% =====================================================================
%  ICROS 2026 (제어·로봇·시스템학회 논문지) — 공식 국문양식(260108) 정밀 반영
%  불법 수상정 포획을 위한 강화학습 기반 무인수상정 협력 제어
%  ▶ 컴파일: XeLaTeX 또는 LuaLaTeX (kotex + fontspec)
%  ▶ 리뷰 1차 대응 반영(구조 재편 / 각도 (sin,cos) / Gerstner 명칭 / 차단망 물리모델 /
%     재현성 파라미터 표 / k값 / heading·formation 근거 / LOS 축소 / 명칭·표기 통일).
%  ▶ 값은 two-USV-system 브랜치 실제 코드 기준. 일부 실험 항목(ablation·다중시드·Fig11 스케일)은 후순위.
%  ▶ 그림: figures/<파일> 있으면 사용, 없으면 placeholder. 수식 일부는 코드 기준 재구성.
% =====================================================================
\documentclass[9pt,twocolumn,a4paper]{extarticle}

\usepackage{kotex}
\IfFontExistsTF{New Batang}{\setmainhangulfont{New Batang}}{}
\IfFontExistsTF{Times New Roman}{\setmainfont{Times New Roman}}{}
\IfFontExistsTF{Malgun Gothic}{\setsanshangulfont{Malgun Gothic}}{}

\usepackage[a4paper,top=13mm,bottom=13mm,left=20mm,right=20mm,columnsep=10mm]{geometry}
\usepackage{amsmath,amssymb,bm}
\usepackage{graphicx}
\usepackage{booktabs}
\usepackage{multirow}
\usepackage{array}
\usepackage{url}
\usepackage{titlesec}
\usepackage{caption}
\usepackage{setspace}
\usepackage{fancyhdr}

\setstretch{1.5}
\setlength{\parindent}{4mm}

\renewcommand{\thesection}{\Roman{section}}
\renewcommand{\thesubsection}{\arabic{subsection}}
\renewcommand{\thesubsubsection}{\arabic{subsection}.\arabic{subsubsection}}
\titleformat{\section}[block]
  {\centering\sffamily\bfseries\fontsize{9}{12}\selectfont}{\thesection.}{0.5em}{}
\titlespacing*{\section}{0pt}{1\baselineskip}{0.4\baselineskip}
\titleformat{\subsection}[block]
  {\sffamily\bfseries\fontsize{9}{12}\selectfont}{\thesubsection.}{0.4em}{}
\titlespacing*{\subsection}{0pt}{0.5\baselineskip}{0.2\baselineskip}

\captionsetup{font=footnotesize,labelformat=empty,justification=raggedright,singlelinecheck=false}
\newcommand{\figcap}[2]{\refstepcounter{figure}%
  \caption*{\textbf{그림~\thefigure.}~#1\\ \textbf{Fig.~\thefigure.}~#2}}
\newcommand{\tabcap}[2]{\refstepcounter{table}%
  \caption*{\textbf{표~\thetable.}~#1\\ \textbf{Table~\thetable.}~#2}}

\graphicspath{{figures/}}
\newcommand{\figph}[1]{%
  \IfFileExists{figures/#1}{\includegraphics[width=\linewidth]{#1}}%
  {\fbox{\parbox[c][0.40\columnwidth][c]{0.9\linewidth}{\centering\ttfamily\footnotesize \detokenize{#1}\par\smallskip (그림 추가 필요)}}}}

\pagestyle{fancy}\fancyhf{}
\fancyhead[L]{\footnotesize Journal of Institute of Control, Robotics and Systems (2026)}
\fancyhead[R]{\footnotesize ISSN: 1976-5622}
\fancyfoot[L]{\footnotesize Copyright\copyright\ ICROS 2026}
\fancyfoot[R]{\footnotesize \thepage}
\renewcommand{\headrulewidth}{0.3pt}\renewcommand{\footrulewidth}{0.3pt}

\providecommand{\doi}[1]{doi:~\url{#1}}

% --- 저자 약력: 사진 25x30mm + 이름 + 약력 + ORCID ---
\newcommand{\biophoto}[1]{%
  \IfFileExists{figures/#1}{\includegraphics[width=25mm,height=30mm]{#1}}%
  {\fbox{\parbox[c][30mm][c]{23mm}{\centering\scriptsize 사진}}}}
\newcommand{\authorbio}[4]{%
  \par\smallskip\noindent
  \begin{minipage}[t]{25mm}\vspace{0pt}\biophoto{#1}\end{minipage}\hspace{2mm}%
  \begin{minipage}[t]{\dimexpr\linewidth-27mm\relax}\vspace{0pt}%
    \footnotesize\textbf{#2}\par #3\par\smallskip \textit{ORCID:}~\texttt{#4}\end{minipage}%
  \par\medskip}

\begin{document}

% ===================== 제목 / 저자 / 초록 (전단) =====================
\twocolumn[{%
\begin{@twocolumnfalse}
\begin{center}
{\setstretch{1.3}
{\sffamily\bfseries\fontsize{16}{20}\selectfont 불법 수상정 포획을 위한 강화학습 기반 무인수상정 협력 제어\par}
\vspace{1mm}
{\fontsize{15.5}{19}\selectfont Multi-Agent Reinforcement Learning for Cooperative Unmanned\\
Surface Vehicle Control in Illegal Vessel Capture\par}
\vspace{3mm}
{\sffamily\fontsize{10}{13}\selectfont 윤 영 준$^{1\dagger}$, 방 현 태$^{1\dagger}$, 하 지 웅$^{1}$, 전 유 라$^{2}$, 신 헌 용$^{3}$, 윤 원 근$^{1*}$\par}
\vspace{1mm}
{\fontsize{9}{12}\selectfont (Youngjun Yun$^{1}$, Hyuntae Bang$^{1}$, Jiwoong Ha$^{1}$, Youra Jun$^{2}$, Hunyong Shin$^{3}$, and Wonkeun Youn$^{1,*}$)\par}
\vspace{1mm}
{\fontsize{8}{10}\selectfont $^{1}$Department of Autonomous Vehicle System Engineering, Chungnam National University;\\
$^{2}$Department of Mechanical Engineering, KAIST;\quad $^{3}$Agency for Defense Development (ADD)\par}
}
\end{center}
\vspace{2mm}

{\fontsize{9}{13.5}\selectfont \leftskip=7mm \rightskip=7mm
\noindent\textbf{Abstract:} This paper proposes a cooperative capture system in which two allied
Unmanned Surface Vehicles (USVs) jointly deploy a capture net to intercept non-cooperative vessels in
naval environments. To overcome the limitations of traditional rule-based methods under environmental
uncertainties, a Multi-Agent Reinforcement Learning (MARL) framework is developed within a
high-fidelity Unity-based simulation incorporating Gerstner wave and wind disturbance models. Each
agent utilizes an ego-centric polar coordinate observation space; distance values are normalized via a
non-linear rational function to prioritize close-range sensitivity, while angular discontinuities are
resolved through a $(\sin,\cos)$ trigonometric transformation. We adopt the Multi-Agent POsthumous
Credit Assignment (MA-POCA) algorithm within a Centralized Training with Decentralized Execution
(CTDE) architecture, enabling robust strategy learning through posthumous credit assignment even when
agents exit the environment. The reward structure combines sparse event rewards for capture success
with dense continuous rewards based on formation maintenance, target approach, and heading alignment.
Quantitative results show that the proposed approach significantly outperforms modified Line-of-Sight
(LOS) guidance, achieving a 65\% capture success rate and a three-fold reduction in formation collapse
under dynamic sea conditions.\par
\smallskip
\noindent\textbf{Keywords:} unmanned surface vehicles (USVs), multi-agent reinforcement learning (MARL),
cooperative capture, maritime simulation\par
}
\vspace{4mm}
\noindent\rule{\textwidth}{0.3pt}
{\scriptsize\par\smallskip
Manuscript received 2026.xx.xx; revised 2026.xx.xx; accepted 2026.xx.xx.\\
$\dagger$These authors contributed equally to this work. $*$Corresponding Author.\\
윤영준: 충남대학교 자율운항시스템공학과 석사과정 (0204jun@o.cnu.ac.kr).\;
방현태: 충남대학교 자율운항시스템공학과 박사과정 (htbang@g.cnu.ac.kr).\;
하지웅: 충남대학교 자율운항시스템공학과 석사과정 (nb10316@g.cnu.ac.kr).\;
전유라: KAIST 기계공학과 박사과정 (youraj@kaist.ac.kr).\;
신헌용: 국방과학연구소 선임연구원 (newoldd@naver.com).\;
윤원근: 충남대학교 자율운항시스템공학과 조교수 (wkyoun@cnu.ac.kr).\\
※ 이 연구는 2025년 정부(방위사업청)의 재원으로 국방과학연구소의 지원을 받아 수행된 미래도전국방기술 연구개발사업임 (No.~915141301).\par}
\vspace{3mm}
\end{@twocolumnfalse}
}]

% ============================== I. 서론 ==============================
\section{서론}
어선 단속, 해상 밀수 차단, 해양보호구역 침범 선박 제지 등 해상 내 비협조 선박을 저지해야 하는 수요는
민간 해양 경비와 군사 방어 측면에서 중요해졌다. 군사적으로도 자폭 선박과 같은 무인 수상정의 공격이
새로운 위협으로 대두되면서 이에 대한 무력화 방법 연구가 필수적이다~\cite{ref1}. 위협 선박 저지 방식은
하드킬(Hard-Kill)과 소프트킬(Soft-Kill)로 구분된다. 하드킬은 요격 미사일·함포 등으로 직접 파괴하여
제압 확실성을 제공하나 막대한 비용과 살상을 수반한다~\cite{ref2}. 소프트킬은 비살상 무력화 수단으로 GPS
재밍/스푸핑 기반 전자전이 대표적이나, 자체 관성항법(INS)·광학 자율항법 선박에는 효과가 제한적이고 아군
통신까지 교란하는 부수적 피해가 존재한다~\cite{ref3}.

이러한 한계를 극복하고자 본 연구에서는 두 대의 수상정이 차단망을 전개하여 침입 선박의 추진 기관을
물리적으로 구속하는 방식을 제안한다. 이는 전자 시스템에 비의존하여 아군 통신에 영향을 주지 않으면서도,
저비용·비살상이라는 점에서 기존 하드킬·소프트킬의 한계를 동시에 보완한다. 다만 유인 포획은 탑승자 안전을
보장할 수 없어 무인 선박 도입이 요구된다~\cite{ref4}. 이를 위해서는 두 수상정 간 정교한 협력 기동,
차단망의 기하학적 제약, 해양 환경의 불확실성을 복합적으로 고려해야 한다.

전통적 규칙 기반 접근법은 사전 정의된 시나리오에서는 효과적이나, 파도·바람 외란과 위협 선박의 불규칙한
접근 경로 등 실시간 변동 요소에 유연하게 대응하기 어렵다. 반면 심층 강화학습(DRL)은 시행착오 상호작용으로
복잡한 의사결정을 스스로 해결하며 충돌 회피·경로 계획 등 해상 환경에서 효용이 입증되었다~\cite{ref5,ref6}.
특히 다중 에이전트 강화학습(MARL)은 협력 과제에서 에이전트 간 상호작용 전략을 자동 학습하여 군집 대형
유지, 다중 표적 포획, 추적·포위 기동 등에 적용되었다~\cite{ref7,ref8}. 기존 연구는 주로 CTDE 구조로 학습
시 전역 정보를 활용하고 실행 시 분산 제어를 달성하였으나~\cite{ref9}, 표적 포획 문제에 MARL을 적용한 예는
드물다. 본 연구는 MARL 기반 차단망 협력 포획 학습 프레임워크를 제안한다(그림~\ref{fig:loop}).

\begin{figure}[t]\centering
\figph{fig1_loop.png}
\figcap{물리 환경 시뮬레이션을 포함한 에이전트 상호작용 루프}{The Loop of Agent Interaction Within the Physics Simulation Environment}
\label{fig:loop}
\end{figure}

본 논문의 주요 기여는 다음과 같다. 첫째, 파도·바람 모델을 포함한 Unity 기반 해상 시뮬레이션 환경을
구축하여 협력 포획 학습 프레임워크를 제안한다. 둘째, 에이전트 중심 극좌표 기반 관측 체계와 거리 범위·각도
불연속 문제를 해결하는 관측 공간 설계 방법을 제시한다. 셋째, MARL 구조로 두 아군 USV가 차단망을 공유하며
협력 포획 기동을 자율 학습함을 정량적 비교 실험으로 입증한다.

% ============================== II. 포획 시스템 ==============================
\section{포획 시스템}
\begin{figure}[t]\centering
\figph{fig2_capture.png}
\figcap{차단망을 이용한 협력 포획 시스템}{Cooperative Net-based Capture System}
\label{fig:capture}
\end{figure}

제안하는 포획 방식은 그림~\ref{fig:capture}와 같이 두 아군 선박이 편대를 이루어 선박 사이에 전개한 차단망으로
적군 선박을 포획하는 협력 기동이다. 차단망은 두 선박의 실시간 위치에 따라 동적으로 형상이 변하며, 선박 간
거리에 비례하여 유효 전개 폭이 결정된다. 차단망의 유효성은 두 선박 간 거리 $d_s=\lVert \bm{p}_1-\bm{p}_2\rVert$로
결정되며, 활성 조건은 식~(\ref{eq:netactive})와 같다.
\begin{equation}\label{eq:netactive}
\text{net}=
\begin{cases}
\text{active}, & d_{\min}\le d_s \le d_{\max}\\
\text{inactive}, & \text{otherwise}
\end{cases}
\end{equation}
여기서 $d_{\min}=20$\,m는 차단망이 충분한 장력을 확보하는 최소 활성화 거리이며, $d_s<d_{\min}$이면 그물이
접혀 포획력을 잃는다. $d_{\max}=120$\,m는 협력 대형이 유지되는 최대 간격으로, 이를 초과하면 대형 붕괴로
간주한다. 활성 상태에서 적군 선박이 두 선박을 잇는 그물 선분과 접촉하면 포획이 성립한다(식~\ref{eq:capture}).
\begin{equation}\label{eq:capture}
\text{capture} \iff (\text{net}=\text{active}) \;\wedge\; \big(\mathrm{dist}(\bm{p}_e,\,\overline{\bm{p}_1\bm{p}_2})\le \tfrac{w}{2}\big)
\end{equation}
$\bm{p}_e$는 적 위치, $\overline{\bm{p}_1\bm{p}_2}$는 그물 선분, $w$는 그물 폭이다. 이는 실제 해상에서 차단망이
일정 수준 이상 전개되어야만 유효한 포획이 가능하다는 물리적 제약을 반영한다.

% ============================== III. 시뮬레이션 환경 및 보상 설계 ==============================
\section{시뮬레이션 환경 및 보상 설계}
\subsection{훈련 환경}
강화학습 환경을 위해 Unity 엔진으로 해상 시뮬레이터를 구축한다. 시뮬레이터는 방어 대상인 모선, 이를
방어하는 아군 선박 2대, 모선을 향해 접근하는 적군 선박으로 구성된다(그림~\ref{fig:sim}). 본 연구에서는 자폭
선박의 돌진 상황을 학습 환경으로 설정하였다. 적군은 약 1{,}000\,m 거리에서 스폰되어 모선을 향해 최단 경로로
고속 접근하며(접근 속도는 매 에피소드 $\pm30\%$ 무작위), 모선 반경 20\,m 이내 진입 시 방어 실패(돌파)로
처리된다. 아군 2대는 제한 시간 내 협력 기동으로 적군을 차단·포획해야 한다.

\begin{figure}[t]\centering
\figph{fig3_sim.png}
\figcap{Unity 기반 해상 시뮬레이션 구현 결과: (a) 협력 기동 두 아군 USV, (b) 방어 실패 시 모선 피격, (c) 차단망으로 연결된 두 아군 USV·적군 USV}{Implementation of the Unity-Based Maritime Simulation}
\label{fig:sim}
\end{figure}

\subsection{외란 구현}
해양 물리는 Gerstner 파랑 모델로 구현하였다~\cite{ref10,ref11}. 다중 파형 합성으로 해수면 변위를 계산하며,
수직·수평 변위는 식~(\ref{eq:gv}),~(\ref{eq:gh})와 같다.
\begin{equation}\label{eq:gv}
y(\bm{X},t)=\sum_{i=1}^{N} A_i \sin\theta_i
\end{equation}
\begin{equation}\label{eq:gh}
\begin{bmatrix}x\\ z\end{bmatrix}=\sum_{i=1}^{N} Q_i A_i\,\bm{D}_i\cos\theta_i
\end{equation}
$A_i$는 진폭, $\bm{D}_i$는 단위 진행 방향, $Q_i$는 첨예도 계수이다. 각 파형의 파장 $\lambda_i$로부터
\textbf{파수(wavenumber)} $k_i=2\pi/\lambda_i$가 정의되고(식~\ref{eq:k}), 심해 분산관계로부터
\textbf{각진동수} $\omega_i$(식~\ref{eq:omega})와 \textbf{위상 속도} $c_i$(식~\ref{eq:c})가 결정된다.
\begin{equation}\label{eq:k}
k_i = \frac{2\pi}{\lambda_i}
\end{equation}
\begin{equation}\label{eq:omega}
\omega_i = \sqrt{g\,k_i}
\end{equation}
\begin{equation}\label{eq:c}
c_i = \frac{\omega_i}{k_i}
\end{equation}
위상각은 $\theta_i = k_i(\bm{D}_i\cdot\bm{X}) - \omega_i t$이며, 마루의 날카로움(첨예도)은
$Q_i=Q/(\omega_i A_i N)$로 제어한다($Q$=첨예도 상수, $N$=총 파형 수). 선체 풍력은 매 물리 스텝
$\bm{F}_{\text{wind}}=s\,V_{\text{wind}}\,\hat{\bm{u}}_{\text{wind}}$로 외력 적용한다. 매 에피소드 파도·풍력을
무작위로 설정하는 도메인 랜덤화로 외란 강건성과 sim-to-real 전이 가능성을 높였다~\cite{ref12}.

\subsection{보상 설계}
보상체계는 기하 기반 연속 보상과 임무 결과 기반 순간 보상을 계층적으로 결합한다.

\textbf{대형 유지 보상}은 두 선박 간격 $d_p$가 차단망 최적 전개 간격 $d^{*}=100$\,m에 가까울수록 커지는
가우시안 형태이다(식~\ref{eq:rform}). 표준편차 $\sigma=40$\,m는 차단망 길이 한계 내 허용 분산으로,
$|d_p-d^{*}|=\sigma$에서 약 0.37로 감쇠한다. 즉 단순 임계가 아니라 최적 간격을 중심으로 부드럽게 보상하여
그물이 과도하게 접히거나($d_p<d_{\min}$) 붕괴되지($d_p>d_{\max}$) 않도록 유도한다.
\begin{equation}\label{eq:rform}
R_{\text{form}}=\alpha_f\,\exp\!\Big(-\frac{(d_p-d^{*})^2}{2\sigma^2}\Big)
\end{equation}
\textbf{적 접근 보상}은 거리 감소량이 양수일 때만 부여하여 적군을 향한 실질 접근을 유도한다(식~\ref{eq:rapp}).
\begin{equation}\label{eq:rapp}
R_{\text{app}}=\alpha_a\,\max(0,\; d_{t-1}-d_t)
\end{equation}
\textbf{선수 정렬 보상}은 진행 방향과 적군 방향의 코사인이 양수일 때 부여한다(식~\ref{eq:rhead}). 차단망 포획에서
핵심은 적의 진행 경로를 가로막는 \emph{배치}이며, 이는 대형 유지·접근 보상이 담당한다. 선수 정렬은 접근·차단
기동을 보조하는 정렬 신호로서 가장 작은 가중치($\alpha_h=0.0003$)를 부여하여, 배치 목표와 충돌하지 않고
보완적으로만 작용하도록 설계하였다.
\begin{equation}\label{eq:rhead}
R_{\text{head}}=\alpha_h\,\max(0,\;\cos\theta)
\end{equation}
\textbf{시간 벌점}은 매 스텝 $-0.0001$로 신속한 포획을 유도한다. \textbf{순간(이벤트) 보상}은 포획 성공
$+1.0$(포획 쌍 개별 $+0.3$), 모선 피격 $-1.0$, 아군 충돌 $-0.5$, 전멸 보너스 $+2.0$, 종료 시 잔존 적
1척당 $-0.5$이다. 포획·모선피격 시 에피소드는 조기 종료되고, 충돌 시 해당 에이전트는 비활성화된다. 총 보상은
식~(\ref{eq:rtot})와 같다~\cite{ref18,ref19}.
\begin{equation}\label{eq:rtot}
R=R_{\text{form}}+R_{\text{app}}+R_{\text{head}}+R_{\text{time}}+R_{\text{event}}
\end{equation}
재현성을 위한 보상·환경·정규화 주요 파라미터는 표~\ref{tab:param}에 정리하였다.

\begin{table}[t]\centering
\tabcap{주요 보상·환경·정규화 파라미터}{Key reward, environment, and normalization parameters}
\label{tab:param}
\smallskip
\begin{tabular}{|l|l|c|}
\hline
구분 & 파라미터 & 값\\ \hline
\multirow{4}{*}{보상 가중치}
 & 대형 유지 $\alpha_f$ & 0.002\\ \cline{2-3}
 & 적 접근 $\alpha_a$ & 0.001\\ \cline{2-3}
 & 선수 정렬 $\alpha_h$ & 0.0003\\ \cline{2-3}
 & 시간 벌점 & $-0.0001$\\ \hline
\multirow{5}{*}{이벤트 보상}
 & 포획 (+개별) & $+1.0\ (+0.3)$\\ \cline{2-3}
 & 모선 피격 & $-1.0$\\ \cline{2-3}
 & 아군 충돌 & $-0.5$\\ \cline{2-3}
 & 전멸 보너스 & $+2.0$\\ \cline{2-3}
 & 잔존 적 (1척) & $-0.5$\\ \hline
\multirow{3}{*}{대형·차단망}
 & 최적 간격 $d^{*}$ & 100 m\\ \cline{2-3}
 & 허용 분산 $\sigma$ & 40 m\\ \cline{2-3}
 & 활성 간격 $[d_{\min},d_{\max}]$ & $[20,120]$ m\\ \hline
\multirow{3}{*}{정규화 $k$}
 & 적 거리 $k_e$ & 77\\ \cline{2-3}
 & 파트너 거리 $k_p$ & 100\\ \cline{2-3}
 & 방위각 $k_\beta$ & 10\\ \hline
\multirow{3}{*}{환경}
 & 적 스폰 거리 & 1000 m\\ \cline{2-3}
 & 모선 돌파 반경 & 20 m\\ \cline{2-3}
 & 적 속도 무작위 & $\pm30\%$\\ \hline
\end{tabular}
\end{table}

% ============================== IV. 강화학습 문제 정식화 ==============================
\section{강화학습 문제 정식화}
\subsection{학습 구조: MA-POCA와 CTDE}
두 선박의 협력 포획에서는 한 선박이 포획을 완료하거나, 아군 충돌·대형 붕괴로 \emph{도중에 비활성화}되는 일이
빈번하다. MADDPG(Multi-Agent Deep Deterministic Policy Gradient)~\cite{ref14}나 QMIX~\cite{ref15}는 고정된
에이전트 수를 전제하므로, 이러한 에이전트 이탈 상황에서 학습이 불안정해진다. 본 연구는 사후 크레딧 할당
(Posthumous Credit Assignment)을 제공하는 MA-POCA~\cite{ref13}를 채택하였다. MA-POCA는 에이전트가 이탈한
이후에도 그 기여를 평가할 수 있어, 포획 완료·충돌 비활성화·대형 붕괴 등 \emph{팀 구성원 수가 동적으로 변하는}
본 포획 시나리오에 직접 부합한다(그림~\ref{fig:mapoca}).

\begin{figure}[t]\centering
\figph{fig4_mapoca.png}
\figcap{MA-POCA 기반 협동 학습 구조~\cite{ref13}}{MA-POCA-based Cooperative Learning Architecture}
\label{fig:mapoca}
\end{figure}

학습 시 중앙 비평가는 전체 $N$개 에이전트 관측을 Attention 계층으로 처리해 팀 가치 $V(\bm{s})$를 추정한다
(식~\ref{eq:value}). 이는 각 선박의 상태를 동적으로 가중·통합하여 협력적 가치를 산출하는 역할을 한다.
\begin{equation}\label{eq:value}
V(\bm{s})=g\!\left(\textstyle\sum_{i=1}^{N}\bm{c}_i\right)
\end{equation}
Attention 가중치(식~\ref{eq:attn})는 에이전트 $i$가 다른 선박 $j$에 두는 집중도를, 맥락 벡터(식~\ref{eq:ctx})는
그 가중 통합 결과를 나타낸다. $N\times N$ 행렬은 에이전트 수가 변해도 자동 조정되어 가변 구성에 대응한다
(그림~\ref{fig:credit}).
\begin{equation}\label{eq:attn}
\alpha_{ij}=\mathrm{softmax}_j\!\left(\frac{\bm{q}_i^{\top}\bm{k}_j}{\sqrt{d_k}}\right)
\end{equation}
\begin{equation}\label{eq:ctx}
\bm{c}_i=\sum_{j=1}^{N}\alpha_{ij}\,\bm{v}_j
\end{equation}

\begin{figure}[t]\centering
\figph{fig5_credit.png}
\figcap{Attention 기반 크레딧 할당 구조}{Attention-Based Credit Assignment Architecture}
\label{fig:credit}
\end{figure}

에이전트 $j$를 제외한 기준선 가치(식~\ref{eq:base})를 이용해, 팀 보상에서 $j$의 순수 기여인 상대적 이득
(식~\ref{eq:adv})을 분리 추정한다. 이는 특정 선박이 이탈해도 나머지 선박의 포획 기여 학습이 지속되도록 한다.
\begin{equation}\label{eq:base}
b_j(\bm{s}) = V\!\left(\bm{s}_{\setminus j}\right)
\end{equation}
\begin{equation}\label{eq:adv}
A_j^{t}=r_t+\gamma\,V(\bm{s}_{t+1})-b_j(\bm{s}_t)
\end{equation}
각 정책 $\pi_\theta$는 이 상대적 이득으로 PPO Clipped Surrogate 목적(식~\ref{eq:ppo})을 최대화하며, 갱신 비율은
식~(\ref{eq:ratio})와 같다~\cite{ref18}. 클리핑은 과도한 정책 변화를 억제하여 학습을 안정화한다.
\begin{equation}\label{eq:ppo}
L^{\text{CLIP}}(\theta)=\hat{\mathbb{E}}_t\!\big[\min(\rho_t \hat{A}_j,\,
\mathrm{clip}(\rho_t,1-\epsilon,1+\epsilon)\hat{A}_j)\big]
\end{equation}
\begin{equation}\label{eq:ratio}
\rho_t=\frac{\pi_\theta(a_t^{j}\mid o_t^{j})}{\pi_{\theta_{\text{old}}}(a_t^{j}\mid o_t^{j})}
\end{equation}
실행 시 각 에이전트는 자신의 관측 $o^{j}$만으로 독립 행동하므로 중앙 통신 없이 분산 실행이 가능하다(CTDE).

\subsection{정책 네트워크}
개별 정책 네트워크는 그림~\ref{fig:policy}와 같이 관측 $o^{j}$를 2개의 완전연결층(FC Encoder)으로 인코딩한 뒤
LSTM 순환 신경망으로 처리한다~\cite{ref16}. 포획은 적의 경로 예측·속도 추정 등 단일 시점만으로 어려운
시계열적 판단을 요구하므로 LSTM을 적용하였다. 게이트는 식~(\ref{eq:lstm1})--(\ref{eq:lstm6})과 같다.
\begin{align}
f_t &=\sigma(W_f[h_{t-1},x_t]+b_f) \label{eq:lstm1}\\
i_t &=\sigma(W_i[h_{t-1},x_t]+b_i) \label{eq:lstm2}\\
\tilde{c}_t &=\tanh(W_c[h_{t-1},x_t]+b_c) \label{eq:lstm3}\\
c_t &=f_t\odot c_{t-1}+i_t\odot \tilde{c}_t \label{eq:lstm4}\\
o_t &=\sigma(W_o[h_{t-1},x_t]+b_o) \label{eq:lstm5}\\
h_t &=o_t\odot\tanh(c_t) \label{eq:lstm6}
\end{align}
은닉 상태 $h_t$는 추력·조타의 2차원 연속 행동을 결정한다. 정책 네트워크 구성은 표~\ref{tab:net}과 같다.

\begin{figure}[t]\centering
\figph{fig6_policy.png}
\figcap{개별 에이전트 정책 네트워크 구조}{Individual Agent Policy Network Architecture}
\label{fig:policy}
\end{figure}

\begin{table}[t]\centering
\tabcap{개별 에이전트 정책 네트워크 구성}{Individual agent policy network configuration}
\label{tab:net}
\smallskip
\begin{tabular}{|l|c|}
\hline
Parameter & Value\\ \hline
Hidden units & 256\\ \hline
Num layers & 2\\ \hline
Memory size & 256\\ \hline
Sequence length & 64\\ \hline
\end{tabular}
\end{table}

\subsection{관측 공간}
각 에이전트는 NED 좌표계 기준 자기 중심(ego-centric) 극좌표로 환경을 관측한다. 절대 좌표는 위치에 따라 관측이
달라져 일반화가 저하되므로, 모든 관측을 자기 위치·헤딩 기준 상대 극좌표로 변환하여 회전 불변성을 확보한다
(그림~\ref{fig:obs}). 관측 대상은 파트너·적군·모선이며, 상태 벡터는 아군 관측 $\bm{o}_{\text{ally}}$와 적군
관측 $\bm{o}_{\text{enemy}}$를 결합한 $\bm{o}=[\bm{o}_{\text{ally}},\bm{o}_{\text{enemy}}]$이다. 각 대상에 대해
거리·상대 방위각·상대 선수각을 관측하며 원시 범위는 표~\ref{tab:obs}와 같다.

\begin{figure}[t]\centering
\figph{fig7_obs.png}
\figcap{에이전트 중심 극좌표 관측 구조}{Agent-Centric Polar Coordinate Observation Framework}
\label{fig:obs}
\end{figure}

\begin{table}[t]\centering
\tabcap{에이전트 관측 공간 구성}{Agent observation space configuration}
\label{tab:obs}
\smallskip
\begin{tabular}{|l|l|c|}
\hline
Target & Observation & Raw Range\\ \hline
\multirow{3}{*}{Partner ship}
 & Distance          & $1$--$50$ m\\ \cline{2-3}
 & Relative bearing  & $-180^\circ$--$180^\circ$\\ \cline{2-3}
 & Relative heading  & $-180^\circ$--$180^\circ$\\ \hline
\multirow{3}{*}{Enemy ship}
 & Distance          & $0$--$5000$ m\\ \cline{2-3}
 & Relative bearing  & $-180^\circ$--$180^\circ$\\ \cline{2-3}
 & Relative heading  & $-180^\circ$--$180^\circ$\\ \hline
\end{tabular}
\end{table}

\subsection{관측 정규화}
파트너 거리(최대 50\,m)와 적군 거리(최대 5{,}000\,m)는 범위 차가 커, 신경망이 절댓값이 큰 적군 거리에 편향될
수 있다. 이를 해결하기 위해 모든 거리를 식~(\ref{eq:nd})의 유계 유리함수로 $[-1,1)$로 정규화하여 근거리
변화율을 키웠다.
\begin{equation}\label{eq:nd}
\bar{d}=\frac{d}{\,|d|+k\,}
\end{equation}
상수 $k$는 출력이 0.5가 되는 기준 거리로, 관측 대상별로 달리 설정한다(적군 $k_e=77$, 파트너 $k_p=100$,
방위각 $k_\beta=10$). 적군은 원거리에서도 위협을 감지하도록, 파트너는 근접 협력 간격을 민감하게 보도록 조정한
것이다(그림~\ref{fig:dsens}).

\begin{figure}[t]\centering
\figph{fig8_dsens.png}
\figcap{관측 대상별 거리 감도 조정}{Distance Sensitivity Tuning by Observation Target}
\label{fig:dsens}
\end{figure}

각도 관측은 $\pm180^\circ$ 경계에서 불연속이 발생한다. 사인 함수 단독 변환은 $30^\circ$와 $150^\circ$가 동일한
값을 가져 방위의 일대일 대응이 깨진다. 이를 해결하기 위해 본 연구는 식~(\ref{eq:na})와 같이 \textbf{사인·코사인
쌍}으로 변환하여, 모든 각도를 연속이면서 유일하게 표현하고 방향 정보를 온전히 보존한다.
\begin{equation}\label{eq:na}
g(\theta)=\big(\sin\theta,\;\cos\theta\big)
\end{equation}
이는 적군과의 정밀한 각도 정렬이 요구되는 최종 차단 기동에서, 정렬 부근의 미세 편차를 모호함 없이 감지하게
한다(그림~\ref{fig:asens}).

\begin{figure}[t]\centering
\figph{fig9_asens.png}
\figcap{각도 관측값의 삼각함수 정규화 전·후 비교}{Comparison of angular observation before and after trigonometric normalization}
\label{fig:asens}
\end{figure}

% ============================== V. 학습 결과 ==============================
\section{학습 결과}
\subsection{학습 과정}
학습 효율을 위해 한 실행 내 다수의 환경 인스턴스를 동시 구동하는 병렬 방식을 적용하였다(그림~\ref{fig:parallel}).
MA-POCA 학습 하이퍼파라미터는 표~\ref{tab:hyp}과 같다. 학습 결과(그림~\ref{fig:reward}--\ref{fig:vloss})
팀 보상은 약 1.1M 스텝 이후 수렴하며, 엔트로피는 점차 감소해 결정론적 정책으로 수렴한다. 정책·가치 손실의
동반 감소는 MA-POCA의 안정적 수렴을 뒷받침한다. 여기서 그림~\ref{fig:reward}의 보상은 \emph{팀(그룹) 누적 보상}
(Group Cumulative Reward)이다. MA-POCA는 협력 성과를 그룹 단위로 부여하므로 개별 에이전트의 누적 보상은
$\approx 0$이 되고, 팀 전체의 그룹 보상이 실질 학습 신호가 된다.

\begin{figure}[t]\centering
\figph{fig10_parallel.png}
\figcap{다중 환경 인스턴스를 활용한 병렬 학습 구조}{Parallel Training Architecture with Multiple Environment Instances}
\label{fig:parallel}
\end{figure}

\begin{table}[t]\centering
\tabcap{MA-POCA 학습 하이퍼파라미터}{MA-POCA training hyperparameters}
\label{tab:hyp}
\smallskip
\begin{tabular}{|l|c|}
\hline
Parameter & Value\\ \hline
Batch size & 2048\\ \hline
Buffer size & 20480\\ \hline
Learning rate & $1.13\times10^{-4}$ (linear)\\ \hline
Beta (entropy) & $1.15\times10^{-3}$\\ \hline
Epsilon (clip) & 0.2\\ \hline
Lambda (GAE~\cite{ref20}) & 0.95\\ \hline
Num epoch & 3\\ \hline
Gamma & 0.99\\ \hline
Time horizon & 128\\ \hline
Time scale & 10\\ \hline
\end{tabular}
\par\smallskip
{\footnotesize 학습 모델은 약 $1.7\times10^{6}$ 스텝 체크포인트(수렴 $\approx1.1\times10^{6}$ 스텝)를 사용. 학습률·beta는 실제 실행값.}
\end{table}

\begin{figure}[t]\centering\figph{fig11_reward.png}
\figcap{팀(그룹) 누적 보상 수렴 그래프}{Group Cumulative Reward Convergence graph}\label{fig:reward}\end{figure}
\begin{figure}[t]\centering\figph{fig12_entropy.png}
\figcap{엔트로피 그래프}{Entropy graph}\label{fig:entropy}\end{figure}
\begin{figure}[t]\centering\figph{fig13_ploss.png}
\figcap{정책 손실 수렴 그래프}{Policy Loss Convergence graph}\label{fig:ploss}\end{figure}
\begin{figure}[t]\centering\figph{fig14_vloss.png}
\figcap{가치 손실 수렴 그래프}{Value Loss Convergence graph}\label{fig:vloss}\end{figure}

\subsection{성능 비교}
제안 정책을 baseline인 수정 LOS(Line-of-Sight) Guidance와 동일 환경에서 비교하였다~\cite{ref21}. 비교 용이성을
위해 공격선 1대를 가정한다. 동적 환경에서는 고정 경유점 추종이 부적합하므로, 모선과 포획 대상의 실시간 상대
좌표로 가상 경로를 생성하고($\hat{\bm{d}}$, 법선 $\bm{n}$, 목표 $\bm{p}_i^{\text{tgt}}=\bm{p}_e+L_a\hat{\bm{d}}+s_i
L_s\bm{n}$), PD 조향($K_p{=}3,K_d{=}0.5$)과 거리 비례 추력으로 추종하는 수정 LOS를 비교군으로 사용하였다
(그림~\ref{fig:loscmp}).

\begin{figure}[t]\centering
\figph{fig15_los.png}
\figcap{포획을 위한 적응형 LOS Guidance}{Modified LOS Guidance for Target Capture}
\label{fig:loscmp}
\end{figure}

두 정책 모두 동일 조건에서 각 40회 에피소드를 수행하였다(그림~\ref{fig:traj_los},~\ref{fig:traj_rl}). 평가 지표는
포획률(Capture rate), 방어선 돌파율(Penetration rate), 대형 붕괴율(Formation collapse rate, 아군 충돌 또는
간격이 차단망 길이를 초과한 비율)이며, 결과는 표~\ref{tab:result} 및 그림~\ref{fig:bar}와 같다.

\begin{figure}[t]\centering\figph{fig16_traj_los.png}
\figcap{Modified LOS 기반 방어 궤적 및 포획 결과}{Defense Trajectories and Capture Results Based on Modified LOS Guidance}\label{fig:traj_los}\end{figure}
\begin{figure}[t]\centering\figph{fig17_traj_rl.png}
\figcap{강화학습 모델 기반 방어 궤적 및 포획 결과}{Defense Trajectories and Capture Results Based on RL Model Inference}\label{fig:traj_rl}\end{figure}

\begin{table}[t]\centering
\tabcap{제어 방식별 방어 성능 비교}{Defense Performance Comparison by Control Method}
\label{tab:result}
\smallskip
\begin{tabular}{|l|c|c|}
\hline
Evaluation Metric & Modified LOS & RL (proposed)\\ \hline
Capture rate            & 6/40 (15.0\%)  & \textbf{26/40 (65.0\%)}\\ \hline
Penetration rate        & 25/40 (62.5\%) & \textbf{11/40 (27.5\%)}\\ \hline
Formation collapse rate & 9/40 (22.5\%)  & \textbf{3/40 (7.5\%)}\\ \hline
\end{tabular}
\end{table}

\begin{figure}[t]\centering
\figph{fig18_bar.png}
\figcap{제어 방식별 방어 성능 비교(막대그래프)}{Bar Chart of Defense Performance by Control Method}
\label{fig:bar}
\end{figure}

표~\ref{tab:result}에서 강화학습 정책은 모든 지표에서 향상을 보였다. 포획률은 15.0\%에서 65.0\%로, 돌파율은
62.5\%에서 27.5\%로, 대형 붕괴율은 22.5\%에서 7.5\%(약 3배 감소)로 개선되었다. 사전 정의 경로 추종 제어가
동적 적군 기동에 능동 대응하기 어려운 반면, 강화학습 정책은 실시간 최적 행동으로 복합 방어 목표를 동시
달성함을 의미한다. 다만 본 비교는 단일 시드·40회 평가이며 LOS 게인이 고정값이므로, 다중 시드 통계와 MAPPO
등과의 ablation을 통한 정량 검증은 향후 과제로 남긴다.

% ============================== VI. 결론 ==============================
\section{결론}
본 연구는 두 척의 아군 USV가 차단망을 공유하여 불법 선박을 협력 포획하는 문제에 대해 MA-POCA 기반 다중
에이전트 강화학습 접근법을 제안하였다. 외란 해상 시뮬레이션 환경을 구축하고, 차단망의 물리적 활성·포획 조건을
정식화하였으며, 극좌표 관측 체계와 정규화(거리 유리함수, 각도 사인·코사인 쌍)로 관측 공간을 설계하였다. 가우시안
대형 유지·접근·선수 정렬의 연속 보상과 포획·방어 실패의 순간 보상을 계층 결합하여, 수정 LOS Guidance 대비
포획률·대형 붕괴율을 대폭 개선하였다. 향후에는 MAPPO 등과의 ablation 및 다중 시드 통계 검증, 다수 적군·복수
아군 쌍의 가변 분산 군집 제어, 침입 선박의 독립 회피 정책 학습(상호 경쟁적 진화)으로 확장하고자 한다.

% ============================== 참고문헌 ==============================
{\fontsize{9}{13.5}\selectfont
\begin{thebibliography}{99}
\bibitem{ref1} A. Boretti, ``Unmanned surface vehicles for naval warfare and maritime security,'' \textit{J. Defense Modeling and Simulation}, 2024. \doi{https://doi.org/10.1177/15485129241283056}
\bibitem{ref2} National Research Council, \textit{Integration of Hard-Kill and Soft-Kill Systems for More Effective Fleet Air Defense}. Washington, DC: The National Academies Press, 1992.
\bibitem{ref3} S. Tashakori, M. Ranjbar, S. Balochian, J. Sharif-Razavian, and M. Peymankar, ``Dynamic soft-kill weapon-target assignment in naval environments,'' \textit{Computers \& Industrial Engineering}, vol.~197, p.~110606, 2024. \doi{https://doi.org/10.1016/j.cie.2024.110606}
\bibitem{ref4} H. Bang, D. Choi, J. Lee, E.-T. Kim, and W. Youn, ``Development and Experimental Validation of an Autonomous USV Berthing System with Transition Logic,'' \textit{Ocean Engineering}, 2025. \doi{https://doi.org/10.1016/j.oceaneng.2025.121333}
\bibitem{ref5} R. S. Sutton and A. G. Barto, \textit{Reinforcement Learning: An Introduction}, 2nd ed. Cambridge, MA: MIT Press, 2018.
\bibitem{ref6} J. Woo and N. Kim, ``Collision avoidance for an unmanned surface vehicle using deep reinforcement learning,'' \textit{Ocean Engineering}, vol.~199, p.~107001, 2020. \doi{https://doi.org/10.1016/j.oceaneng.2020.107001}
\bibitem{ref7} C. Yu et al., ``The Surprising Effectiveness of PPO in Cooperative Multi-Agent Games,'' \textit{NeurIPS Datasets and Benchmarks Track}, 2022.
\bibitem{ref8} Y. Zhao, Y. Ma, and S. Hu, ``USV formation and path-following control via deep reinforcement learning with random braking,'' \textit{IEEE Trans. Neural Netw. Learn. Syst.}, vol.~32, no.~12, pp.~5468--5478, 2021. \doi{https://doi.org/10.1109/TNNLS.2021.3068762}
\bibitem{ref9} F. A. Oliehoek and C. Amato, \textit{A Concise Introduction to Decentralized POMDPs}. Cham: Springer, 2016. \doi{https://doi.org/10.1007/978-3-319-28929-8}
\bibitem{ref10} A. Lindell, ``Boat Attack,'' GitHub repository, Unity Technologies, 2020.
\bibitem{ref11} J. Tessendorf, ``Simulating Ocean Water,'' \textit{SIGGRAPH 2001 Course Notes}, ACM, pp.~1--26, 2001.
\bibitem{ref12} J. Tobin, R. Fong, A. Ray, J. Schneider, W. Zaremba, and P. Abbeel, ``Domain Randomization for Transferring Deep Neural Networks from Simulation to the Real World,'' \textit{IEEE/RSJ IROS}, pp.~23--30, 2017. \doi{https://doi.org/10.1109/IROS.2017.8202133}
\bibitem{ref13} A. Cohen et al., ``On the Use and Misuse of Absorbing States in Multi-Agent Reinforcement Learning,'' arXiv:2111.05992, 2021.
\bibitem{ref14} R. Lowe, Y. Wu, A. Tamar, J. Harb, P. Abbeel, and I. Mordatch, ``Multi-Agent Actor-Critic for Mixed Cooperative-Competitive Environments,'' \textit{NeurIPS}, pp.~6382--6393, 2017.
\bibitem{ref15} T. Rashid et al., ``QMIX: Monotonic Value Function Factorisation for Deep Multi-Agent Reinforcement Learning,'' \textit{ICML}, pp.~4295--4304, 2018.
\bibitem{ref16} S. Hochreiter and J. Schmidhuber, ``Long Short-Term Memory,'' \textit{Neural Computation}, vol.~9, no.~8, pp.~1735--1780, 1997. \doi{https://doi.org/10.1162/neco.1997.9.8.1735}
\bibitem{ref17} A. Vaswani et al., ``Attention Is All You Need,'' \textit{NeurIPS}, pp.~5998--6008, 2017.
\bibitem{ref18} J. Schulman, F. Wolski, P. Dhariwal, A. Radford, and O. Klimov, ``Proximal Policy Optimization Algorithms,'' arXiv:1707.06347, 2017.
\bibitem{ref19} A. Y. Ng, D. Harada, and S. Russell, ``Policy Invariance Under Reward Transformations,'' \textit{ICML}, pp.~278--287, 1999.
\bibitem{ref20} J. Schulman, P. Moritz, S. Levine, M. I. Jordan, and P. Abbeel, ``High-Dimensional Continuous Control Using Generalized Advantage Estimation,'' \textit{ICLR}, 2016.
\bibitem{ref21} T. I. Fossen, M. Breivik, and R. Skjetne, ``Line-of-sight path following of underactuated marine craft,'' \textit{IFAC Proc. Volumes}, vol.~36, no.~21, pp.~211--216, 2003. \doi{https://doi.org/10.1016/S1474-6670(17)37809-6}
\end{thebibliography}
}

% ===================== 저자 약력 =====================
\section*{저 자 소 개}
\authorbio{bio_youngjun.png}{윤 영 준}{2025년 충남대학교 자율운항시스템공학과 학사 졸업. 2026년 3월$\sim$현재 충남대학교 자율운항시스템공학과 대학원 석사과정. 관심분야는 강화학습 기반 경로 생성 및 무인 수상정 군집 제어.}{0009-0005-8321-0461}
\authorbio{bio_hyuntae.png}{방 현 태}{2019년 한밭대학교 건설환경공학과 학사 졸업. 2021년 한밭대학교 토목공학과 석사 졸업. 2022년$\sim$현재 충남대학교 자율운항시스템공학과 대학원 박사과정. 관심분야는 무인 이동체 제어 시스템 및 검증 시뮬레이션.}{0000-0002-2846-6334}
\authorbio{bio_jiwoong.png}{하 지 웅}{2024년 한남대학교 전자공학과 학사 졸업. 2025년$\sim$현재 충남대학교 자율운항시스템공학과 대학원 석사과정. 관심분야는 무인 이동체 제어 시스템 및 HILS 시스템 구축.}{0009-0005-2980-5413}
\authorbio{bio_youra.png}{전 유 라}{2021년 아리조나대학교 기계공학과 석사 졸업. 2025년 9월$\sim$현재 한국과학기술원(KAIST) 기계공학과 대학원 박사과정. 관심분야는 유한요소 해석 기반 구조 해석 및 최적화.}{0000-0002-9591-2533}
\authorbio{bio_hunyong.png}{신 헌 용}{2007년 부산대학교 전자전기통신공학과 학사 졸업. 2009년 광주과학기술원 정보통신공학과 석사 졸업. 2012년$\sim$2021년 국방기술품질원 선임연구원. 2021년$\sim$현재 국방과학연구소 선임연구원. 관심분야는 전투체계, 심층강화학습, 무인수상정 경로계획.}{0009-0006-7247-4215}
\authorbio{bio_wonkeun.png}{윤 원 근}{2021년 KAIST 전기및전자공학부(로봇 학제전공) 박사 졸업. 2021년 3월$\sim$현재 충남대학교 자율운항시스템공학과 조교수. 관심분야는 무인이동체 유도$\cdot$제어$\cdot$항법 및 응용, HILS 시스템 개발 및 검증.}{0000-0001-6954-9059}

\end{document}
